\section{An actual nuclear endomorphism on the nontrivial character theta ideal}
\label{sec:character-common-sum-transport}

In a nontrivial finite-order compact character sector the source polar
representation is absent. This removes the two output conditions of
MA-MO-010, but it does not by itself prove the required inverse domains.
We first prove an unconditional division statement for every actual
smooth theta-ideal vector. We then identify the complete original
common-sum action on that core and extend its smoothing to its actual
Hilbert closure. This is an endomorphism of the completed theta ideal;
it is not a construction of the full smooth zero quotient or its trace.

Let \(K\) be a fixed number field and let \(\eta\) be a nontrivial
primitive finite-order Hecke character. This includes a nontrivial
primitive Dirichlet character and the non-induced cubic character
\(\psi_{269,3}\) already constructed over \(\mathbb Q(\zeta_3)\).
Choose a structural finite place set \(S\) as in MA-CC-001,
containing the conductor and the required additive-character data.
Its finite part is understood in Euler products. Use the same set
for \(\overline\eta\). Put
\[
 \begin{aligned}
 A_{\omega,S}(s)&=E_{\omega,S}(s)L_\omega(s),\\
 E_{\omega,S}(s)&=\prod_{\mathfrak p\in S}
          (1-\omega(\mathfrak p)(N\mathfrak p)^{-s}),\\
 \mathscr I^c_\omega&=T_{1/2}\mathcal I_\omega\\
       &=\{y\in\mathcal E:(\mathcal My)^{(j)}(\rho-1/2)=0\\
       &\hspace{5em}\text{for }\rho\in Z_\omega,
                                      \ 0\le j<m_\rho\},\\
 H_{\Sigma,a,s}&=L^2\bigl(\mathbb R,
         (e^{2ax}+e^{-2ax})(1+x^2)^{s/2}\,dx\bigr),\\
 J_{\omega,a,s}&=\overline{\mathscr I^c_\omega}^{\,H_{\Sigma,a,s}},\\
 &0<a<1/2,\quad s\ge0,\quad\omega\in\{\eta,\overline\eta\}.
 \end{aligned}
 \tag{MA-MO.170}
\]
Here \(Z_\omega\) is the full nontrivial completed zero divisor,
with multiplicities. The equality defining \(\mathscr I^c_\omega\) is the exact
MT-M18/MT-M24 image theorem, not a newly defined Euler ideal.
All constants concern the fixed field, character, place set and width.
Use the probability compact Haar measure and norm splitting of MT-M24
and MA-CC.8. Thus \(C_g\) is the actual centered scalar convolution
associated with the compact-character test in MA-CC.5.

Let \(\phi\) be the original real source cutoff, equal to one
sufficiently far left and zero sufficiently far right. It need not
be symmetric. We use either of the following two regimes:
\[
 \begin{array}{ll}
 \text{gap/smooth:}&
  |1/2+a-\operatorname{Re}\rho|\ge\delta>0
     \quad(\rho\in Z_\eta\cup Z_{\overline\eta}),\\
 &g,\phi'\in C_c^\infty;\\[2mm]
 \text{generic/Gevrey:}&
  1/2+a\in\mathscr G_{\eta,\overline\eta},\\
 &g,\phi'\text{ are compact Gevrey of a common finite order }q>1.
 \end{array}
 \tag{MA-MO.171}
\]
Here \(\mathscr G_{\eta,\overline\eta}\) is the simultaneous
full-measure set of MA-RS-005. The gap is an explicit hypothesis.
The generic regime has no
zero-location premise.

\paragraph{MA-MO-012 (actual character common-sum smoothing on its ideal closure).}
For every \(y\in\mathscr I^c_\eta\), the complete original
source action of \(C_g(P_+^S-P_-)\) on \((y,y)\) is a diagonal
pair \((B^{\rm raw}_{\eta,g}y,B^{\rm raw}_{\eta,g}y)\), with
\(B^{\rm raw}_{\eta,g}y\in\mathscr I^c_\eta\).
In either regime (MA-MO.171), this action extends uniquely to
an actual trace-class endomorphism
\[
                   B_{\eta,g}:J_{\eta,a,s}\longrightarrow J_{\eta,a,s}
                                      \qquad(s\ge0).
 \tag{MA-MO.172}
\]
The map is continuous in trace norm on every fixed compact smooth
test-support space in the gap regime, and uniformly bounded on
families with fixed support and common finite-order Gevrey bounds
in the generic regime. No additional moment conditions on \(g\)
or on the source vector are imposed.

The proof includes the following domain statement without either
analytic regime: for every \(y\in\mathscr I^c_\omega\), its actual
original wide inverse \(F=J^{\rm geom}_{\omega,S}y\) satisfies
\[
 \begin{gathered}
 e^{bx}F(x)\in\mathscr S(\mathbb R)\quad\hbox{for every }b>-1/2,
       \qquad (1-\phi)F\in\mathcal E,\\
 J_{b,\omega,S}y=F\quad(b=\pm a),\\
 Z_{b,\omega,S}(\phi F)=Z^{\rm geom}_{\omega,S}(\phi F)
                                         \in\mathscr I^c_\omega.
 \end{gathered}
 \tag{MA-MO.173}
\]
The narrow symbols in this formula denote the maximal multipliers
on the indicated lines. The vector \(F\) need not belong to
\(\mathcal E\); the weighted statement is the one proved and used.

\begin{proof}
\emph{Unconditional division on the positive finite-\(L\) half-plane.}
For \(y\in\mathscr I^c_\omega\), set
\[
             H_y(z)=\mathcal My(z-1/2),\qquad
                         Q_y(z)=H_y(z)/L_\omega(z).
 \tag{MA-MO.174}
\]
The numerator is entire and rapidly decreasing, with all
derivatives, on every fixed closed vertical strip. All nontrivial
completed zero jets vanish to their full order. Since the gamma
factors have no zeros, the possible poles of \(Q_y\) that these
jets do not remove occur only at the trivial finite-\(L\) zeros.
For a finite-order character those occur among the nonpositive
real gamma-pole locations. There are only finitely many in any
fixed vertical strip. This also covers a possible overlap at
an endpoint: its remaining pole order is at most the gamma-pole
order, after the full completed-zero multiplicity is canceled.

Fix a central closed strip \(\epsilon\le\Re z\le R\), with
\(\epsilon>0\). Choose a wider strip \(-A\le\Re z\le B\),
where \(A>2\), \(B>\max(2,R)\), and the left edge avoids the
finitely located trivial-zero real parts. Let \(P(z)\) be a
finite polynomial clearing all the possible trivial-zero poles
of \(Q_y\) in that wider strip; its zeros can all be taken in
\(\Re z\le0\). The orders can be fixed from the gamma factors,
independently of \(y\). Then
\[
                         P(z)Q_y(z)
       \quad\hbox{is holomorphic on the wider closed strip.}
 \tag{MA-MO.175}
\]

The selected heights of MT-M18, (M18.7)--(M18.9), give polynomial
lower bounds for the finite \(L_\omega\), uniformly on each
horizontal segment of this wider strip. Indeed the exponential
Stirling factor in the lower bound for the completion cancels
against the same factor in \(\Xi_\omega/L_\omega\).
On the two chosen vertical edges the finite function likewise
has a polynomial lower bound: at the right edge use the Euler
product; at the left use the completed functional equation,
Stirling, and the right-half-plane Euler lower bound. The left
edge was chosen to avoid the finite low-height gamma exceptions.
All these statements are unconditional for the fixed character.

For each sufficiently large positive or negative height block,
the rapid bounds for \(H_y\), multiplied by \(P\) and these
polynomial reciprocal bounds, therefore give, for every \(N\),
\[
       |P(z)Q_y(z)|\le C_{N,y}(2+|n|)^{-N}
 \quad\hbox{on the boundary of }
                  [-A,B]+i[h_n,h_{n+1}].
 \tag{MA-MO.176}
\]
The maximum principle gives the same estimate inside each
rectangle. The heights have gaps between \(1/2\) and \(3/2\),
so these estimates cover both infinite ends of the central strip.
On that strip \(P\) is bounded away from zero at bounded height
and has polynomial size at large height. The finitely many
remaining portions of the central strip are compact and \(Q_y\)
is holomorphic there. Thus \(Q_y\) is rapidly decreasing on that strip.
Cauchy estimates on a slightly wider positive strip give the
same conclusion for every derivative.

Only now divide by the finite deleted Euler product. On
\(\Re z\ge\epsilon>0\), every nontrivial factor obeys
\(|1-\omega(\mathfrak p)(N\mathfrak p)^{-z}|
\ge1-(N\mathfrak p)^{-\epsilon}>0\), and its inverse and
all derivatives are uniformly bounded on each such strip.
We have proved
\[
 \widetilde Q_y(z)=\frac{H_y(z)}{E_{\omega,S}(z)L_\omega(z)},
 \qquad \widetilde Q_y\in\operatorname{Hol}(\{\Re z>0\}).
 \tag{MA-MO.177}
\]
This function is rapid with all derivatives on every positive closed
strip. No lower bound for \(E_{\omega,S}\) was used on a strip crossing
\(\Re z=0\), where its zeros need not be finite in number.

\emph{The actual inverse and the cutoff domain.}
Inverse Mellin integration of
\(\widetilde Q_y(z+1/2)\) along any \(\Re z=b>-1/2\)
defines the same smooth function \(F\). The contours may be
shifted between any two such lines by (MA-MO.177), with no
poles and vanishing horizontal integrals. Fourier inversion
gives
\[
 \begin{gathered}
 \mathcal MF(z)=\frac{\mathcal My(z)}
                   {E_{\omega,S}(z+1/2)L_\omega(z+1/2)}
                  \quad(\Re z>-1/2),\\
             e^{bx}F\in\mathscr S\quad(b>-1/2).
 \end{gathered}
 \tag{MA-MO.178}
\]
The Mellin identity is understood through the holomorphic
quotient, including its canceled zero values. On every wider
line \(b>1/2\), it is exactly the original absolutely
convergent Euler inverse of MA-CC.1. Hence this \(F\) is
\(J^{\rm geom}_{\omega,S}y\), the source inverse, rather
than a chosen lift of the Hilbert vector. The source
identification is valid for the structural set \(S\) by
MA-CC.8 and Meyer's Lemmas~5.3--5.4.
More explicitly, MT-M24 gives an actual positive-image class for
\(y\). Its compact character is invariant outside this structural
set \(S\), so Meyer's local identification supplies a vector of
the corresponding \(\mathcal T(\mathcal C_S)\) sector whose
summation is \(y\). The wide Euler inverse is unique, and
therefore that vector is exactly \(F\). This verifies its local
source membership before any source Fourier map is applied.

The weights with arbitrarily large positive \(b\) show
rapid decay of \(F\) and all its derivatives at the right
end. The cutoff \(q=1-\phi\) vanishes at the left end.
Consequently
\[
                           qF\in\mathcal E.
 \tag{MA-MO.179}
\]
For any \(U\in\mathcal E\), the literal ideal sum
\(Z^{\rm geom}_{\omega,S}U\) belongs to
\(\mathscr I^c_\omega\). This uses only the nontrivial
finite character, not universality: on a wide line its
Mellin transform is \(A_{\omega,S}(z+1/2)\mathcal MU(z)\).
It is entire with rapid decay on every fixed strip, because
\(L_\omega\) is entire and has the standard polynomial
strip bounds. Shifting to every line gives an element of
\(\mathcal E\); its full completed-zero jets vanish, so
MT-M24 identifies it with the actual theta image. This is
the residue-free part of MA-MO-009, which applies here
without the degree-two universality hypothesis.

Apply it to \(U=qF\). Linearity of the actual wide Euler
sum gives the literal identity
\[
       Q^{\rm geom}_\omega y
          :=Z^{\rm geom}_{\omega,S}(\phi F)
           =y-Z^{\rm geom}_{\omega,S}(qF)
                                   \in\mathscr I^c_\omega.
 \tag{MA-MO.180}
\]
All these original sums exist on the wide source domain.
For \(b=\pm a\), write
\(A_{\omega,b}(t)=A_{\omega,S}(1/2+b+it)\).
Formula (MA-MO.178) gives
\(V_by=A_{\omega,b}V_bF\), with \(V_bF\) Schwartz.
Multiplication by \(A_{\omega,b}^{-1}\) therefore returns
\(V_bF\) almost everywhere, even if the line has zeros.
This verifies the maximal inverse domain. Both \(F\) and
\(qF\) are in the maximal \(Z_{b,\omega,S}\) domain, so
their difference \(\phi F\) is too. Its output is precisely
(MA-MO.180). This proves all of (MA-MO.173) without a gap
or a generic-width assumption.

\emph{The actual complete common-sum formula and the dual sector.}
A finite-order character is trivial on the connected norm
splitting. Since \(\eta\ne1\), its compact restriction is
nontrivial. Both source polar characters are trivial on that
compact group, so their averages in this sector vanish.
The source Poisson identity therefore makes the entire
positive image diagonal; by MT-M24 it is exactly
\(\mathscr I^c_\eta\). These are statements about the
specified compact average and image, not a radial model
assumed for a single ideal class.

For \(y\in\mathscr I^c_\eta\), let \(F\) be its actual
local source inverse and put
\(\widetilde F=J^{\rm geom}_{\overline\eta,S}Ry\).
The actual source embeddings and approximate section give
\[
 \begin{aligned}
 i_+F&=(y,y),\qquad
            \widetilde F=\mathfrak F^*F,\\
 \bigl(P_+^S(y,y)\bigr)_0
       &=Z^{\rm geom}_{\eta,S}(qF)
          +Z^{\rm geom}_{\eta,S}\mathfrak F(q\widetilde F),\\
 Z^{\rm geom}_{\eta,S}\mathfrak F(q\widetilde F)
       &=RZ^{\rm geom}_{\overline\eta,S}(q\widetilde F).
 \end{aligned}
 \tag{MA-MO.181}
\]
Here \(\mathfrak F,\mathfrak F^*\) are Meyer's actual
additive Fourier maps on the local coinvariant spaces,
including the finite-place phases. They are not replaced
by a bare real or complex archimedean Fourier transform.
The last identity is the source Poisson formula, whose
singular term is zero in this nontrivial compact sector.
The cutoff vectors \(qF,q\widetilde F\) belong to
\(\mathcal E\) by the division statement and are valid
local source vectors. These assertions also respect the
odd Dirichlet Fourier-square phases already checked in
MA-CQ; no separate assumption of an archimedean involution
is used.

Inversion in MA-CC.9 exchanges \(\eta\) and
\(\overline\eta\) and becomes the linear reflection \(R\).
In particular
\[
                 R\mathscr I^c_\eta=\mathscr I^c_{\overline\eta}.
 \tag{MA-MO.182}
\]
The primitive finite local factors and the conductor remain
those of \(\eta\); reflection conjugates their coefficients.
Using \(Z^{\rm geom}_{\eta,S}F=y\) and
\(Z^{\rm geom}_{\overline\eta,S}\widetilde F=Ry\) in
(MA-MO.181), and subtracting \(P_-(y,y)=(y,y)\), gives
\[
 \begin{aligned}
 (P_+^S-P_-)(y,y)&=(\Delta_\eta y,\Delta_\eta y),\\
 \Delta_\eta y&=y-Q^{\rm geom}_\eta y
                         -R Q^{\rm geom}_{\overline\eta}Ry
                                      \in\mathscr I^c_\eta.
 \end{aligned}
 \tag{MA-MO.183}
\]
The second component agrees by the zero-singular-term
common-sum relation. Every term here is an actual source
function; the preceding argument checked all its inverse
and cutoff domains. Convolution preserves the actual
theta image by equivariance, or directly by multiplication
of all its Mellin jets. Thus
\(B^{\rm raw}_{\eta,g}y=C_g\Delta_\eta y\) belongs to
\(\mathscr I^c_\eta\).

\emph{Trace-class kernels on both lines and both character sectors.}
The completed functional equation sends a zero \(\rho\) of
\(\eta\) to \(1-\rho\) of \(\overline\eta\), with the
same multiplicity. Hence the union of the two divisors is
stable under this reflection. The gap or generic-distance
bound at \(1/2+a\) supplies the same bound at \(1/2-a\).
The disk minimum-modulus proofs of MA-RS-005 apply to
both lines: their positive real parts lie inside the same
fixed inner disk at large height, and there are no trivial
zeros there. The reciprocal derivatives are polynomial
in the gap regime and bounded by
\(\exp(C_j\log^2(2+|t|))\) in the generic regime.
Finite Euler inverses have bounded derivatives on these
positive closed strips. No estimate for a reciprocal on
the negative half-plane is invoked here.

For \(\omega=\eta,\overline\eta\) and \(b=\pm a\), put
\(j_{\omega,b}=A_{\omega,b}^{-1}\),
\(G_b(t)=\mathcal Mg(b+it)\), and
\(P(w)=\int e^{iwx}\phi'(x)\,dx\). The sequential
core difference \(D_{\omega,b}=M_\phi-
Z_{b,\omega,S}M_\phi J_{b,\omega,S}\), after smoothing,
has the exact Fourier kernel
\[
 k_{\omega,b,g}(t,u)
   =\frac{i}{2\pi}G_b(t)A_{\omega,b}(t)
            \frac{j_{\omega,b}(t)-j_{\omega,b}(u)}{t-u}P(t-u).
 \tag{MA-MO.184}
\]
The canceled diagonal has its ordinary limiting value.
The RS proofs show that this is a Schwartz kernel: in the
gap regime use polynomial reciprocal derivatives and
Schwartz decay; in the generic regime the two Gevrey
Fourier tails dominate the logarithmic-square growth.
The argument controls every mixed derivative and applies
to all four line/character choices.

Write \(K_{\omega,b}(g)\) for its extension. It is trace
class on every \(H_{b,s}=L^2(e^{2bx}(1+x^2)^{s/2}dx)\).
For completeness, with \(r=s/2\) and
\(\Lambda^r=(1-\partial_t^2)^{r/2}\),
\[
 \begin{gathered}
       \|u\|_{H_{b,s}}=\|\Lambda^rV_bu\|_2,\\
 k^{(r)}_{\omega,b,g}
                  =\Lambda_t^r\Lambda_u^{-r}k_{\omega,b,g}.
 \end{gathered}
 \tag{MA-MO.185}
\]
The second line is the kernel of the operator conjugated by
\(\Lambda^r\) in the \(V_b\) Fourier picture. It remains
Schwartz, since the two Bessel multipliers and their
derivatives have polynomial growth. Its Hermite expansion
is absolutely summable in rank-one trace norms. This proves
trace class in the actual weighted space, as in MA-MO-010.

\emph{Agreement of these kernels on every smooth ideal vector.}
The unconditional division lemma supplies
\(u=V_by\in\mathscr S\) and
\(q=j_{\omega,b}u=V_bF\in\mathscr S\) for every
\(y\in\mathscr I^c_\omega\). In either regime the line
has no zeros. Choose compact smooth frequency cutoffs
\(\chi_N\to1\), and set \(u_N=\chi_Nu\). Then
\[
       u_N\longrightarrow u,\qquad
       j_{\omega,b}u_N=\chi_Nq\longrightarrow q
                              \quad\hbox{in Schwartz topology}.
 \tag{MA-MO.186}
\]
Multiplication by \(\phi\) in physical coordinates and
by \(A_{\omega,b}\) in Fourier coordinates preserves
Schwartz topology. Thus the sequential outputs converge
there, and in every Sobolev norm. The bounded kernel
extensions agree on \(u_N\) and hence on \(u\).
This proves agreement with the actual geometric outputs
(MA-MO.180). It does not assume that the reciprocal
multiplier sends every Schwartz function to Schwartz
in the generic regime. The same argument applies to
\(Ry\) in the dual sector.

\emph{Assembly in the inherited sum norm.}
Let \(\psi(x)=1-\phi(x)-\phi(-x)\), a compact smooth
function, and \(g^\vee(x)=g(-x)\). For each \(b=\pm a\)
the actual smoothing in (MA-MO.183) agrees on its entire
smooth ideal core with
\[
        T_{\eta,b}(g)=C_gM_\psi+K_{\eta,b}(g)
                                +RK_{\overline\eta,-b}(g^\vee)R.
 \tag{MA-MO.187}
\]
Indeed \(D_{\eta,b}+RD_{\overline\eta,-b}R\) has
the multiplication part \(\phi+\phi(-\cdot)\), while
(MA-MO.183) has the multiplication part one.
Reflection is unitary \(H_{b,s}\to H_{-b,s}\) and
\(C_gR=RC_{g^\vee}\). The remaining term \(C_gM_\psi\)
has the Schwartz Fourier kernel
\(G_b(t)\widehat\psi_+(t-u)/(2\pi)\), so is also
trace class in every polynomial weight. Thus (MA-MO.187)
is trace class on the specified one-sided space.

Set \(\mathbb H=H_{a,s}\oplus H_{-a,s}\), and define
\[
 \iota:H_{\Sigma,a,s}\longrightarrow\mathbb H,
       \qquad\iota y=(y,y),\qquad
 \mathbb T_{\eta,g}
       =\operatorname{diag}(T_{\eta,a}(g),T_{\eta,-a}(g)).
 \tag{MA-MO.188}
\]
The map \(\iota\) is isometric with closed range.
On every \(y\in\mathscr I^c_\eta\), the two components
of \(\mathbb T_{\eta,g}\iota y\) agree with the same
actual source vector \(B^{\rm raw}_{\eta,g}y\), which
belongs to \(\mathscr I^c_\eta\). Passing to the defined
closure therefore gives a diagonal output in
\(\iota J_{\eta,a,s}\) for every input from that subspace.
Consequently
\[
 \begin{aligned}
 B_{\eta,g}&=\left.\iota^*\mathbb T_{\eta,g}\iota
                                           \right|_{J_{\eta,a,s}},\\
 \|B_{\eta,g}\|_1
       &\le\|T_{\eta,a}(g)\|_{1,H_{a,s}}
                         +\|T_{\eta,-a}(g)\|_{1,H_{-a,s}}.
 \end{aligned}
 \tag{MA-MO.189}
\]
On this domain \(\iota^*\) is the inverse diagonal
identification of the already compatible outputs; it
does not average inconsistent extensions. Composition
of trace-class and bounded maps gives (MA-MO.172).
The range lies in the same closed ideal, so this time
the map is an endomorphism. Density gives uniqueness.
The finite Schwartz-seminorm estimates in the kernel
proof give all the stated test-family bounds.
\end{proof}

\paragraph{The legitimate trace and the projection that cannot be omitted.}
Let \(\iota_J=\iota|_{J_{\eta,a,s}}\) and
\(\Pi_J=\iota_J\iota_J^*\), the orthogonal projection
onto its actual diagonal ideal in \(\mathbb H\).
Bounded trace cyclicity gives
\[
       \operatorname{Tr}_{J_{\eta,a,s}}B_{\eta,g}
          =\operatorname{Tr}_{\mathbb H}(\mathbb T_{\eta,g}\Pi_J).
 \tag{MA-MO.190}
\]
For example, one absolutely summable rank-one representation
of \(\mathbb T_{\eta,g}\) gives both sides as the same
series, with the same \(\Pi_J\) pairing. This is an
actual operator identity, not a prescription by zero sums.
It does not permit dropping \(\Pi_J\), and hence does
not identify this trace with the sum of the two unrestricted
one-sided traces or with Meyer's full common-space trace.
The endomorphism has its ordinary Fredholm determinant,
but no arithmetic evaluation of it is asserted here.

\paragraph{Requested characters, source placement, and the remaining gate.}
For a nontrivial primitive Dirichlet character, the stipulated
input A supplies the gap in (MA-MO.171) at every fixed width.
For a non-induced imaginary-quadratic character, its own
Hecke RH, or the stipulated input M, supplies that gap;
the generic/Gevrey assertion for almost every width is
unconditional. In particular the result applies to both
\(\psi_{269,3}\) and its conjugate with their primitive
ramified factors and structural place set from MA-CH/MA-CQ.
The principal character is excluded: its polar terms lead
to the distinct domain and correction in MA-MO-007--011.

The exact inherited sources are MT-M18 (M18.4) and the
complete selected-height and minimum-modulus proof
(M18.7)--(M18.9); MT-M24 (M24.3), (M24.6)--(M24.9) for
the actual averaged theta image; Meyer (2005),
\texttt{lem:summation}, \texttt{lem:Hilm\_p\_embed}, the
defining \(P_+^S,P_-\) maps and
\texttt{eq:Poisson\_summation}, pp.~39--42; and MA-CC.8--.9
for the centered compact-character and dual-sector factors.
The inverse proof uses those full multiplicity and actual
image statements, not a guessed adelic representative.

The finite-order character is the previously specified
rank-one finite-monodromy local system; its Euler factor
is the corresponding Frobenius determinant. The analytic
norm-half centering here is not an integral Tate twist.
In the inherited cote~1 trace framework, scans~079--080
and~111--112 through BR-TR-01/02, (MA-MO.189) supplies an
actual nuclear endomorphism and (MA-MO.190) its valid
pairing. No weight-zero positivity or global zero bound
is inferred from that placement.

Finally the domain is the completed \emph{theta ideal},
not all of \(\mathcal E\) completed in the narrow norm.
The proof does not establish the original action on every
smooth zero-quotient class, identify the full Meyer trace
with a narrow quotient trace, remove the projector in
(MA-MO.190), or erase the discarded-divisor gate of
MA-FT-002. Those assertions require additional maps or
estimates. The improvement here is the actual original
common-sum endomorphism on every smooth ideal input and
its trace-class extension in the explicitly stated regimes.
